
% \section{Unbounded combinatorics with confined canonical WAD}
\section{Pull-off Argument: unbounded combinatorics}

\subsection{Setting}\label{assumptions sec} 
%\subsection{Assumptions}\label{assumptions sec} 

Consider the class $\QQ\equiv \QQ(\ul \qq, \bar\bb)$
of renormalizable pseudo-ql maps $f: U' \ra U $ with the elephant eye
combinatorics with parameters $ \qq\equiv \qq_f \geq  \ul\qq  $ and $\bb\equiv \bb_f\leq \bar \bb $,
and  renormalization periods $p\equiv p_f= \qq+\bb$.
Let $\AAA=\AAA_f$ stand for  the Canonical WAD of $f$  in $U\sm \KK$. 

Take an integer  $ N \in (1,     \qq-\bb]  $.
Let $\Om\subset U$ be  a topological disk 
in the translation region 
% with $\di \Om \subset U\sm \KK$
which is the union of 
$2N+1$ consecutive  sectors $\Sec_n\supset K_n$,
$n \in \II = [l-N, l+N] \subset \Z/ p\Z$.
Note that $\Om$ contains the chain of $2N$ short arcs
connecting the $K_n$, $n\in \II$
(or rather: appropriate representatives of these arcs). 
In what follows, we will also use notation $\De_m$ for $\Sec_m$. 

\comm{******
\ssk\nin (T1) 
$$
    \Om= \bigcup_{m=- N }^N \De_m,\quad \di \De_m\subset U\sm \KK,
$$
where the $\De_m$ are Jordan disks with disjoint interiors such that $f$ univalently maps $\De_m$ onto $\De_{m+1}$;
each $\De_m$ contains a block $\BB_m\subset \II$ of  $d$ little Julia sets, i.e.,
$$    
                        \De_m\cap \KK = \bigcup_{n\in\BB_m } K_n; 
$$ 
 

%\ssk \nin (T1) There is a ``central block'' $\BB \subset \II$ of $d$ little Julia sets such that
%$$
%    \II = \bigsqcup_{m=-N}^N \BB_m, \quad \{\mathrm{where}}\ \BB_m=\BB+m;  
%$$

\ssk \nin (T2)  For any $n\in \II$, at least $\WW_n-O(1)$ part of the local weight $\WW_n$ is supported on $\AAA$
    (with an absolute additive constant);

\ssk\nin (T3) Any horizontal arc  $\alpha\in \AAA^\hor$ that begins in $\De_m$ with $|m| <  N$ 
  is confined to $\De_{m-1}\cup \De_m\cup \De_{m+1}$;
%  and it ends in $\BB+m \pm 1$; 

 % \ssk\nin (T4) Any horizontal arc $\alpha$ that lands in $\BB+m$ with $-N <    m \leq N$ has a legitimate pullback $f^*\alpha$
 %    (that, of course,  lands in $\BB + m-1$).  
 ***************} 


\bignote{In what follows, Ref (T2) should be replaced with Corollary \ref{loc weights comparison} (ii) } 
  

 Such a domain $\Om$ will be called a {\it translation domain} of size
 $2N+1$ centered at $K_l$. 
% with  base $\De=\De_0$.
Let us also   impose the following property on such a domain:

\ssk\nin (P1)
All local weights $\WW_n$, $n\in \II$,  are $(1-\eps)$-comparable.

\ssk\nin (P2)
$(1-\eps)$-proportion of the local weight $\WW_n$, $n\in \II$,
is carried by two short arcs emanated from $\KK_n$.

\ssk Here $\eps>0$ can be selected arbitrary small as long as
$\WW_\loc$ and $p$ are big enough.  

\ssk Note that (P1) is easily enforced because all the colca weights are
$2-$comparable (Lemma \ref{comparison of dyn loc w}). 
Property (P2) can be enforced to Corollary \ref{loc weights comparison}. 


 We say that  $\Om\cap \KK$ is the {\it $\KK$-slice} of $\Om$.
% (and similarly for any other domain, for that matter).  
Let  $(f^n)^*\Om $ be the pullback of $\Om$ under $f^n$ centered at $K_{l-n}$.
(We will see later that  actually  the pullback  $(f^n)^*\Om $
% for $n\leq 3p$
is independent of the base.) 

We let
$$
  \Om[\kappa]= \bigcup_{|m| \leq \kappa N} \De_m. 
$$

The Julia sets $\KK_{\pm N} $ will be called {\it peripheral}. Let
$\pp\II= \KK_N\cup \KK_{-N}$ be the family  of two peripheral Julia sets
and $\II^\circ= \II\sm \pp\II$ be the complementary family of non-peripheral Julia sets. 



We let  $\BOM$ stand for the pseudo-disk corresponding to $\Om$, and  $(\Bf^n)^*\BOM$ for the corresponding  pullbacks.

For $s\in \II^\circ$, let $\CC(K_s |\, \II) $   be the conductance of the restricted electric circuit $\AAA|\, \II$ with battery $\{K_s , {\pp}\II\} $. 
Thus,  the potential difference 1 between $K_s$ and ${\pp}\II$ 
creates current  $\CC(\BB |\, \II) $ in the restricted circuit. 
By the Series Law, we have: 

% \ssk\nin (T4)
If $s\in \KK_m$ with $|m|\leq  \kappa N$ for some
$\kappa\in (0,1)$ then
\begin{equation}\label{T4}
  \CC (K_s, \II  ) \leq \frac C{(1-\kappa )N} \, \WW_\loc,
\end{equation}
with an absolute $C$. 
% where $C$ depends only on $d$. 

\bignote{ Below Ref (T4) should be replaced with (\ref{T4}) ! }

It is convenient to play with many translation domains $\Om_l$
simultaneously:


\bsk\nin
{\bf Translation Property}.  \note{OK name?} 
For arbitrary big  $N, L\in \N$ 
% function  $L=L(N)\to \infty$  and an arbitrary big $N$ 
there exists  $\underline{p}=\underline{p}(N,L)$ such that any 
map $f\in \QQ$ with $p \geq \underline{p}$ has  $L$ 
pairwise disjoint translation domains $\Om_l $
centered at some $K_l$
of size $2N+1$. 


\msk
Here the index $l$ runs over some subset $\LL\subset \Z/p\Z$ of size $L$. 
All objects associated with $\Om_l$ will be labeled by $l$ if needed, e.g., 
% $\BB_l$,
$\II_l$ etc. \note{more?}  

\comm{*******
\subsection{Example: Multi-elephant eye with confined  canonical
  WAD} \label{confinement assumption}

Let us consider the multi-elephant eye combinatorics with $r$ eyes. 
Let $\Gamma$ be the almost invariant marking of $U\sm \KK$ by Hubbard arcs.

Given a subfamily  $\JJ\subset \Z/p\Z$ of little Julia sets,  
let us say that an arc $\alpha$ is {\it confined to $\JJ$} if it does not cross 
marked Hubbard edges  attached to any $K_m$ with  $m\not\in \JJ$. 

\bignote{Formulate in terms of sectors}

Let $\JJ_n = \{ n+jq \}_{j=0}^{r-1}$ be the ``fundamental family' 'of Julia sets
containing $K_n$, and let 
$$
   \JJ_n(R) = \bigcup_{ |k |\leq R}  ( \JJ_n +k ) 
$$ 
stand for its ``$R$-neighborhood''. 

For the moment, let us label the little Julia sets from $-p/2$ to $p/2$. 

\msk\nin {\bf Confinement Assumption:}
Any canonical arc that begins on  $K_n$ with $|n| \geq q/N$ 
is confined  to the region  $ \JJ_n(p/N)$. 
% $\{m:\,   |m-n|< p/N\} \subset \Z/p\Z$.

\bignote{Make it stronger: confinement to any region of size $q/N$
  with two buffers of the same size.}

\msk
Under this assumption,
 the canonical WAD well inside the translation region descends to the Ecalle-Voronin cylinder ${\Cyl}$.
Let us refine it to a triangulation of  ${\Cyl}$ with vertices $K_n$ and $\pm i \infty$.
Then we can connect $+i\infty$ to $-i\infty$ with a simple curve $\gamma$ composed of edges  of the triangles. 
 Cutting $\Cyl$ along $\gamma$ we obtain a topological strip in  $\Cyl=\C/\Z$ 
that lifts to a fundamental domain % $\Delta$ 
on $\C$.  
Taking the union of its $2N+1$ 
% forward and backward  
translates, we obtain $\Om$. 

********** } 


\subsection{Amplification Theorem}


\begin{thm}% [Amplification]
\label{amplification thm}
% Under the Translation Assumption,
For any amplification factor $\la > 1$ 
there exist $\underline{\WW}_\loc = \underline{\WW}_\loc(N) $ and  $\underline{p}=\underline{p}(, \la, N,L) $
% (``very  big'')
such that for all sufficiently big $N$ and $L$ % and sufficiently big $W_\loc$ (in terms of $N$),
we have:
$$
     \WW^\ver \geq \la \, \WW_\loc,
$$ 
as long as $\WW_\loc\geq \underline{\WW}_\loc$ and $p\geq \underline{p}$. 
% where $C(N)\to \infty $ as $N\to \infty$.
\end{thm}

\subsection{Definite local loss}

Let us first show that if the local vertical weight is definite in  some pullback $(\Bf^n)^* \BOM$ with $n\leq 3p$,
then the Amplification Theorem holds:

\begin{lem}\label{def vert does it}
% Under the Translation Assumption, 
Assume there exists a  $\de > 0$ 
such that each domain  $\BOM_l$ has a  pullback $(\Bf^n)^*\BOM_l$ with $n\leq 3p$
that  satisfies
$$
      \WW ((\Bf^n)^*\BOM_l\sm K_{s-n}) \geq \de \, \WW_\loc, 
$$
where $K_{s-n}\in (\Bf^n)^* \BOM_l$ and the vertex $K_s$, $s\in\II_l$, 
stays distance at most $\kappa N$ from $K_l$. 
Then there exists  $L=L(N)$ such that the conclusion of the Amplification Theorem holds.
% that for all $N$ big enough we have: 
% $$
%     \WW(\AAA^\ver) \geq \la \, \WW_\loc,
% $$ 
% as long as $\WW_\loc\geq \underline{\WW}_\loc(N)$ and $p\geq \underline{p}(N,L)$. (Here quantification is the same as in  teh Amplification Theorem.}
\end{lem}

% Here the distance is the combinatorial distance in the canonical arc diagram $\AAA$
%%%  (where each arc is assigned length 1).   

\begin{proof}
We let $\BOM\equiv \BOM_l$ % and  $\BOM_{-n}\equiv (\Bf^n)^*\BOM$, 
for simplicity. 

  Let us apply the Covering Lemma to  the map $\Bf^n: (\Bf^n)^* \BOM \sm K_{s-n}\ra \BOM\sm K_s$.
The local degree of this map is at most $8$; the big degree is at most $8^N$. The buffer is given by the maximal immersed annulus
in $\BOM\sm \KK$ around $K_s$; its width, $\WW_s \asymp  \WW_\loc$, is comparable with $\WW ((\Bf^n)^* \BOM \sm K_{s-n})$, by the assumption. 
Hence
$$
      \WW(\BOM\sm K_s) \asymp \WW((\Bf^n)^*\BOM\sm K_{s-n} ) \asymp \WW_\loc. 
$$    


% Let us show that this leads to  amplifications of the vertical weight. Indeed, 
Any vertical path in
the annulus $\BOM\sm K_s$ either crosses  a peripheral
Julia set $K_m$,  $m\in \pp \II$, 
or  overflows a vertical path in $\BOM\sm\KK$ beginning on some  non-peripheral Julia set $K_m$, $m\in \II^\circ$. 

In the former case, our path follows the horizontal  electric circuit
in $\BOM$. By  (\ref{T4}),       % property  (T4),
the width of this
family is bounded by $\frac C {(1-\kappa) N}\, \WW_\loc$ for $N$ big enough.  %   where $\eps(N)\to 0$ as $N\to 0$. 

The width of the latter  path family is bounded by 
$$
  \sum_{m\in \II^\circ} \WW^\ver (K_m|\, \BOM\sm \KK)  %  \sum_{m\in \II^\circ} \WW (K_m\ra \KK \cup \di U  |\, U \sm \KK) +O(N)  
                        = \sum_{m\in \II^\circ} \WW^\ver (K_m) +O(N),  
$$
where the last equality follows from  Corollary \ref{loc weights comparison} (ii).
% property (T2)
% and the definition of the non-peripheral vertices
% (saying that no horizontal canonical arcs connect $K_m$ to $K_n$,
where   $m\in \II^\circ$, $n\not\in  \II$).
% (here we use notation $\II$ of the union of the $K_n$ over $n\in \II$).  
Putting the above remarks together, we conclude that for some $n\in \II^\circ$, 
$$
   \WW^\ver (K_m) \geq \frac c {N} \WW_\loc.
$$
% with  $c>0$ depending only on $d$. 

Let us now recall that by the Translation Property we can do it for $L$ domains $\BOM_l$, $l\in \LL$, with disjoint $\KK$-slices. 
This gives as $L$  different Julia sets $K_{m_l}$,  $m_l\in \II_l $. Letting $L= 2\la N/c $,
we obtain
$$  
   \WW^\ver\geq \sum_{l\in \LL}  (\WW^\ver (K_{m_l})- 2) \geq   \la \WW_\loc. 
$$
\end{proof}


\bignote{Section ``Essential arcs'' is hidden below}

\comm{**********

\subsection{Essential arcs}

For $\xi\in (0,1)$, 
we say that a horizontal arc $\alpha\in \AAA$ is $\xi$-{\it essential} if $\WW(\alpha)\geq \xi \WW_\loc$. 
We let $\AAA^\ess$ be the weighted arc diagram comprised of all essential arcs
(we will often let $\xi$ be implicit in the terminology and notation). 

For a simple closed curve $\gamma\subset U\sm \KK$, 
let the {\it flux} $\WW^\perp (\gamma)$ be the the width of the
annulus associated with $\gamma$ that covers $U\sm \KK$.

\begin{lem}\label{definite flux}
  Assume $p$ is the smallest renormalization period of $f$. 
 Then  for any natural $1\leq d<p $ there exists $\si =\si (d)>0 $ such that for any Jordan curve $\gamma$ surrounding $d$ little Julia sets
we have $\WW^\perp (\gamma) \geq \si \cdot \WW_\loc$.
% wthere $\WW^\perp (\gamma)$ is the flux through $\gamma$. 
\end{lem}

\bignote{Note that ``$d$'' is reserved for the number of Julia sets in the
  fundamental sector.}

\begin{proof}
   Let us do induction in $d$. For $d=1$ we can let $\si=1/2$. Let us pass from $d-1$ to $d$. 
Of course, we can assume that $\si(1)\geq \si(2)\geq\dots \geq \si(d-1)$.

Let $\gamma$ surrounds $d$ little Julia sets $K_n$, $n\in \BB\subset \Z/p\Z$, and let $\De$ be the disc bounded by $\gamma$.

\msk\nin
{\em Claim:} There exists $\xi=\xi(d)$ such that if $\WW^\perp(\gamma)< \xi/2$ then the graph of $\xi$-essential arcs in $\De$
is connected. 

\msk Since $\si(d-1)$ is well defined  inductively, we can let 
$$
   \xi(d)=  \frac {\si(d-1)} {3(d+1)}. 
$$
Assume $\BB=\BB_1\sqcup \BB_2$, where $\BB_1\not=\emptyset$ and $\BB_2\not=\emptyset$ 
cannot be connected with $\xi$-essential arcs in $\De$. Since there are at most $3(d+1)$ 
arcs (vertical and horizontal) in $\De$, and the total weight eminated from $\BB_1$ is at least $\si(d-1)$, 
one of the arcs eminated from $\BB_1$ has weight  at least $\xi (d)$. By the assumption, this arc must be vertical, 
so $\WW^\perp(\gamma) \geq  \xi$ --  contradiction. 

\msk  

We let now 
\begin{equation}\label{si(d)}
  \si(d) =  2^{-(d+2)}\si(d-1) < \xi(d)/2. 
\end{equation}
Assume 
\begin{equation}\label{bound on the flux}
   \WW^\perp(\gamma)< \si (d). 
\end{equation}
Then by the Claim the graph of the $\xi$-essential arcs in $\De$ is connected. 

Select some little Julia set $K_n$ surrounded by $\gamma$, 
and let $\om := (f^p)^\# \gamma$  be the pullback of $\gamma$ based at $K_n$ (immersed into $U\sm \KK$).    
Since $\deg(f^p: \om\ra \gamma) \leq 2^d$, we have  the usual estimate
$$
    \WW^\perp(\om) \leq 2^d \WW^\perp(\gamma) < \si (d-1).
$$
By the induction assumption, $\om$ cannot surround less than  $d$ little Julia sets, 
so it must suround all the $K_n$, $n\in \BB$ (and of course, nothing else). 

Moreover, $\om$ cannot cross any of the $\xi$-essintial arcs in $\De$, for otherwise
$\WW^\perp(\om) \geq \xi$ contradicting (\ref{si(d)}) and (\ref{bound on the flux}). 
Since the graph of these arcs is connected,
$\om$ is homotopic to $\gamma$ in  $U\sm \KK$. But then $f$ is renormalizable with intermediate period
$p/d \in (1,  p) $.  
\end{proof}

% Select some $\xi= \xi(d) > 0$, to be specified below. 
% We say that a horizontal arc $\alpha\in \AAA$ is {\it essential} if $\WW(\alpha)\geq \xi \WW_\loc$. 

\begin{lem}\label{ess chains}
There exist $\xi= \xi(d)$ and  $C=C(d)$ such that for any $m$ with $|m|< N-l $ we have the following property:

\ssk \nin {\rm (i)} Either there is a vertical arc $\alpha$ emanating from $\De_m$  with $\WW(\alpha)> \xi$, 

\ssk \nin {\rm (ii)} Or else,   any $K_n \in \BB_m$  can be connected to 
each  neighboring block $\BB_{m+i}$, $i\in \{\pm 1\}$,
by a path of length $\leq C$ of $\xi$-essential arcs confined to 
$$
         \bigcup_{ |n-m|\leq C} \De_n. 
$$  
\end{lem}

\begin{proof}
Taking $\xi < \si(d)/3(d+1) $, we ensure that there an arc $\alpha$ eminated from $\De_m$ with $\WW(\alpha)> \xi$. 
If this arc is vertical then option (i) holds. Otherwise each $\De_m$ can be connected to a neigboring block, 
$\BB_{m+1}$ or $\BB_{m-1}$, with an essential arc.   % \note{remark on periodicity}

Let us orient the arcs $\alpha\in \AAA^\ess$ so that each $\alpha$ connects $\BB_m$ to $\BB_{m+i}$ with $i\geq 0$. 
  Let us push the $\BB_m$ and $\AAA^\ess$ to the cylinder $\Cyl$; 
we denote the corresponding vertices and arcs downstairs by $\bar K_n$ and  $\bar \alpha$. 
Let us consider the recurrent core $\bar \RR$ of $\bar \BB$ consisting of the vertices $\bar K_n$  that can be included in a loop
of $\bar\AAA^\ess$. Restriction of $\bar \AAA^\ess$ to $\bar\RR$ is ``connected with orientation'': any vertex can be connected to any other
by an orinented path of length bounded in terms of  $d$.
%It is comprised of blocks $\bar \RR_i$ which are cyclically permuted by the graph.  
Hence for any $K_n$ and $K_s$ in $\RR_m\subset \BB_m$ (the lift of
$\bar \RR$ to $\BB_m$)
\note{connectivity issue}
there exists $i\geq 0$ (bounded in terms of $d$) such that 
 $K_n$ can be connected to $K_{s+i}$ in $\AAA^\ess$. 
Let us show that $i>0$ for some of these $i's$. Otherwise, the  connected component of $\RR_m$ (disregarding the orientation)
would be a finite isolated block in $\AAA^\ess$ whose size can be bounded {\it a priori } by some  $d'$ depending only on $d$.
This is impossible provided $\xi< \si(d')/ 3(d'+1)$.  

Since for any $K_t\in \RR_m$,  $K_{s+i}$ can be further connected  to $K_{t+j}$ with some (controlled) $j\geq i$, we see that 
 for any $K_n$ and $K_t$ in $\RR_m$ there exists $i > 0$ (bounded in terms of $d$) such that 
 $K_n$ can be connected to $K_{t+i}$ in $\AAA^\ess$.  Reversing orientation, we obtain the same property with $i<0$. 

Finally, if $\bar K_n\not \in \bar \RR$ then it can be connected to some $\bar K_s\in \bar \RR$ 
by a (disoriented) path of $\bar \AAA^\ess$ of length bounded in terms of $d$. The conclusion follows. 
\end{proof}

******* } 

\bignote{Section ``Loosing Julia sets'' is hidden below}

\comm{****
\subsection{Loosing Julia sets}

Let us say that a set $K_m$, $m\in \II_l$, {\it gets lost} under the $n$-th pullback
if $K_{m-n}\not\subset (f^n)^* \Om_l$. Note that then $K_{m+k}$ gets lost under any further pullback $(f^k)^*$, $k\geq n$. 

\begin{lem}\label{getting lost}
Assume that for each  set $K_l$, $l\in \LL$,  the following property holds:
There exists a vertex $K_j$  of the graph $\AAA|\, \II_l$ on distance
at most $\kappa N$ from $K_l$, 
 which gets lost under $n$-th pullback,   $n \leq  3 p$. Then the Amplification Thereom is valid. 
\end{lem}
 
\begin{proof}
Our assumption implies that  $K_j$ gets lost under the $3p$-th pullback, thus 
\begin{equation}\label{pulled-off J-set}
   K_j \not\subset (f^{3p} )^*\Om , \quad \Om \equiv \Om_l.
\end{equation}    \note{regular $\Om$}
By Lemma \ref{ess chains}, $K_j$ can be connected to $K_l$ by a chain of essential arcs confined to 
\begin{equation}\label{cup BB+m}
   \bigsqcup_{|m| \leq  \kappa  N+C} \BB_m.
\end{equation}
%  Without loss of generality,  
% we can assume that  $K_j$ is the closest to $K_l$ vertex in this chain with property  (\ref{pulled-off J-set}).   
Then there exists $K_s\subset  (f^{3p} )^*\Om $ and 
$K_t\not\subset (f^{3p} )^*\Om $, $|s|, |t|\leq \kappa N+C$,
connected by an essential arc $\alpha$.
But then $\alpha$ restricts to a vertical arc in  $ (\Bf^p)^*\BOM $ of definite weight that begins on $K_s$.
By Lemma \ref{def vert does it}, this implies  the Amplification Thereom.  
\end{proof}

\begin{rem}
  We see, in partucular, that assuming  the contrary to the
  Amplification Theorem, 
the pullback $(f^n)^*\Om$ is independent of the base Julia  set $K_l$. 
\end{rem}


% Let $\{K_n\}_{n\in \II}$ be the family of little Julia sets contained in $\Om_0$,
% $\II\subset \Z/p\Z$. 

\begin{lem}\label{interesting moments}
Assume that for each $l\in \LL$ there exists a pullback $(f^n)^* \Om_l\not\supset K_1$, $n\leq 3p$,   
containing a horizontal arc $\alpha$ with endpoints on distance at
most $\kappa N$ from  $K_{l-n}$
such that  both  lifts  $f^* \alpha$  are pulled off.
Then  the concluison of the  Amplification Theorem  holds. 
\end{lem}
 
\begin{proof}
Take some $l\in \LL$, and let $\Om \equiv \Om_l$. 

Let the arc $\alpha$ lands on $K_i$ and $K_j$. By Lemma \ref{getting lost}, if one of the  Julia sets,  
 $K_{i-1}$ or $K_{j-1}$, gets lost then the concluison of the Amplification Theorem  holds. 
So, assume both $K_{i-1}$ and $K_{j-1}$  are contained in $( f^{n+1})^*\, \Om$. Let $\beta= f^* \alpha$ be the lift of $\alpha$ that begins on $K_{i-1}$. 
If it pulls off then  it ends on $K_{j-1}'$. Hence $K_{j-1}'\subset ( f^{n+1})^*  \Om$. 

Since  $( f^n)^* \Om $ does not contain $K_1$, 
the map $f: ( f^{n+1})^*\, \Om \ra ( f^n)^*\,  \Om $ is univalent. 
It follows that $K_{j-1}\not\subset ( f^{n+1})^*\,  \Om $, so this little Julia set gets lost.
The concluison follows. 
\end{proof}

********************************} 

\bignote{\S ``Good fit'' is hidden below}

\comm{*****
  
\subsection{Good  fit}

We say that a horizontal arc $\alpha$ {\it fits into a domain}
$(f^n)^* \Om$,  
$$ 
     \alpha\hsubset (f^n)^* \Om,
$$ 
 if it is contained in it up to homotopy in $U\sm \KK$ rel $\KK$.  
 If $\alpha\hsubset (f^n)^* (\Om[\kappa]) $,
we say that it  {\em fits $\kappa$-well  into } $(f^n)^* \Om$.

\bignote{Make the terminology and notations consistent!} 

\bignote{For one eye, we have an {\em amazing} fit! (except for a boun
ded number of pullbacks)} 



\begin{lem}\label{fitting of  arcs}
   If an arc $\alpha$ fits $\kappa$-well into some domain
   $(f^n)^*\Om$,  
then it fits $(1-\kappa)$-well into all domains
$(f^{n+j})^*\Om$ for  $j=0,1,\dots,  (1-2\kappa) N $.
\end{lem}

\begin{proof}
      We have: 
$$
          (f^j)^* (\Om[1-\kappa] ) \hsupset \Om [\kappa] , \quad
          0\leq j  \leq      (1-2\kappa)  N. 
$$
Hence 
$$
  (f^{n+j})^* ( \Om [1-\kappa])  \hsupset (f^n)^*(\Om[\kappa]   )
  \hsupset \alpha  ,\quad j=0,1,\dots, (1-2\kappa)N. 
$$ 
\end{proof}

For $\xi\in (0,1)$, 
we say that a horizontal arc $\alpha\in \AAA$ is $\xi$-{\it essential} if $\WW(\alpha)\geq \xi \WW_\loc$. 
We let $\AAA^\ess$ be the weighted arc diagram comprised of all essential arcs
(we will often let $\xi$ be implicit in the terminology and
notation; just  an  apriori selected small number).
              \note{quantify  better?} 

\bignote{This def was transferred from \S ``Essential arcs''.}


\begin{lem}\label{fitting of ess arcs}
  Assume the conclusion of the Amplification Theorem fails.
  Let $\alpha$ be an essential horizontal arc.   
  Then there exists  $l\in \LL$  such that $\alpha$ fits $(1-\kappa)$-well into
   $ (1-2 \kappa) |\II|$ domains $(f^m)^*\Om$, $m< p$,
\end{lem}

\begin{proof}
  There are  at least $ (1-2d\kappa) |\JJ|$ domains  $(f^m)^*(\Om[1-\kappa])$,
  $m< p$, containing one endpoint of $\alpha$. If $\alpha$ does not
  fit into such a  domain, it creates a definite veritcal arc in it.
  If this happens for every $l\in \LL$ then
  by Lemma \ref{def vert does it},   the conclusion of the
  Amplification Theorem is valid.
\end{proof}

**************************}


\subsection{Confinement}

\note{ Move next to Lemma 2.2 ?} 

\begin{lem}\label{localization of local weight}
  Under the Translation Assumption, either the Amplification Theorem holds, or else 
  for any little Julia set $K_n$, there exists  $O(N)$  horizontal arcs $\alpha^n_i$ that begin on $K_n$ 
% and end on  $O(N)$ other little Julia sets $K_m$, 
with total weight
$$
    \sum_i \WW(\alpha^n_i) \geq (1-\eps) \WW_n, 
$$
where $\eps=\eps(N)= O(\sqrt{1/N})$. Moreover, these arcs are confined to the pullback $(\Bf^j)^* \BOM$
of some translation domain $\BOM$
% centered at $K_n$ 
(here $n+j\equiv l \ \mod p$ and $0\leq j< p$). 
\end{lem}

\begin{proof}
Let us consider $L$ translation domains $\BOM_l$, $l\in \LL$. Assume that the vertical weight in each of these domains
emanating from the non-peripheral part $\II_l^\circ$ is at least $\WW_\loc/N$.
Then by  Corollary \ref{loc weights comparison} (ii)  %property (T2)
the vertical weight in 
$U\sm \LL$ emanating from each $\II_l^\circ$ and supported on vertical
arcs is at least $ \WW_\loc/ {2N}$, 
%$\WW_\loc/N - O(N)\geq  \WW_\loc/ {2N}$,
as long as $\WW_\loc$ is sufficiently big (in terms of $N$). 
Thus, if $L > 2 N \la$, we achieve amplification. 

\bignote{The $\Om_l$ should be selected well inside the translation region!} 

So, let us assume that for some translation domain $\BOM\equiv \BOM_l$ we have
$$
  \WW^\ver(\II^\circ|\, \BOM) \leq \frac  {\WW_\loc}N.  
$$
Together with property \ref{T4} this implies:  
$$
   \WW(\BOM\sm K_l) \leq \frac CN \, \WW_\loc. 
$$
By the Covering Lemma,
$$
     \WW( (\Bf^j)^* \BOM  \sm K_n ) \leq  \sqrt{4\,  \WW_\loc \cdot
       \WW(\BOM\sm K_l )} \leq \frac C {\sqrt{N}}\,  \WW_\loc, 
$$
with a constant $C$ depending only on $d$.
Hence the rest of the local weight at $K_n$ is the horizontal weight confined to  $(\Bf^j)^* \BOM$. Since this domain contains 
$O(N)$ little Julia sets, this horizontal weight, up to an additive constant $O(N)$,  is localized on at most $O(N)$ horizontal arcs.
\end{proof}


\subsection{Proof of the Amplification Theorem. }
%
\comm{******
Assume the conclusion of the Amplification Theorem fails.  
If $L$ is big enough, then by Lemmas \ref{getting lost} and \ref{interesting moments},
for some $l\in \LL$ no little Julia set and no horizontal arc (of the type specified in the lemmas) 
get lost (respectively: pulled off) under $3p$ pullbacks of
$\BOM\equiv \BOM_l$;
and moreover, the fitting property of Lemma \ref{fitting of ess arcs} is also satisfied
for this $\Om$.  
Then ******* }
%
Assuming that  the conclusion of the Amplification Theorem fails,  
we will show that one of the domains $(\Bf^n)^*\, \BOM_l$, $ n\leq 3p$, contains a definite vertical weight 
(i.e., of order $N\cdot \WW_\loc$), and the Covering Lemma will complete the proof. 

% Since $l$ is fixed from now on, we will skip it from the notation. 


%Note that to this end,  we can replace $U$ with $U^p$
%(since  we can  push forward a definite vertical arc in $\BOM_{n+p}$ to an definite vertical arc in $\BOM_n$;
%just using that $\deg f^p\leq 2^N$, not the Covering Lemma). So, we  replace $U$ with $U^p$ without changing notation. 

By Lemma \ref{sym vert arcs} (adapted to our situation), the domain 
$V:= ( f^{l+p})^*\, \Om \supset K_0$ \note{!} 
contains two symmetric vertical curves
landing on $K_0$ and disjoint from $\AAA^\hor(V)$. 

Let $\AAA_l$ be the set of horizontal  arcs that begin on $K_l$.  
\note{Defined in \S \ref{assumptions sec}?} 
Then Lemma \ref{in p steps} ensures that all $\alpha\in \AAA_l$
get pulled off in less than $p$ iterates.  
Let $n=n(\alpha)$ be the pull-off moment, i.e.,
there exist a legitimate horizontal pullback $(f^k)^*\alpha $  for $k=0,1,\dots , n-1$, 
while  $(f^n)^* \alpha$ pulls off. 
By  Lemma  \ref{interesting moments},
% \ref{J-sets don't pull off}   and \ref{def vert does it}, 
 $n(\alpha) \in \II_l$. \note{Need subscript $l$?} 


 Given an $\eps>0$, let $\xi=\eps/3d$. 
Let $\AAA_l^\ess$ be the set of $\xi$-essential arcs that begin on
$K_l$.  Their total weight is at least $(1-\eps)\WW_l$ . 

\bignote{Next line: These are the moments when $(f^n)^*\BOM \supset K_1$ !} 

Let $\nu_n$ be the total horizontal  weight in $(\Bf^{n-1})^*\, \BOM'$ that gets pulled off under $\Bf^*$, $n\in \II$,
and let $\eta_n$ be the part of $\WW_0$ that was used at that moment  for creation
of new horizontal arcs. %  compare (\ref{bound by W-1}).
Then the new horizontal weight that can be created in  $(\Bf^n)^*\, \BOM' $ is bounded by 
\begin{equation}\label{bound on nu}
    \frac 14 (2\nu_n  + \eta_n) +O(N). 
\end{equation}
Hence the loss of the horizontal weight as we pass from   $(\Bf^{n-1})^*\, \BOM'$ to  $(\Bf^n)^*\, \BOM'$ is at least
\begin{equation}\label{loss at one step}
    \nu_n  -  \frac 14 (2\nu_n +  \eta_n ) - O(N) = \frac 12 \nu_n - \frac 14 \eta_n  - O(N). 
  \end{equation}\label{nu-prime}

  Let us represent $\eta_n$ as $\sum \eta(\alpha)$
 over the horizontal arcs $\alpha\in (f^{n-1})^* \BOM'  $ that gets
pulled off under $f^*$,
where $\eta (\alpha)$ is the part of $\WW_0$ used for creation of
new horizontal arcs connecting $K_0$ to $f^*\alpha$.
Then we have:
\begin{equation}\label{penalty} 
  \sum \eta_n  = \sum_{n\in \II}  \sum_{\alpha\in (f^{n-1})^*\BOM'}  \eta (\alpha) =
  \sum_\alpha\sum_n  \eta (\alpha) \leq  (2N+1) \sum_\alpha\eta(\alpha) 
  \leq (2N+1) \WW_1. 
\end{equation}  

On the other hand, let us now show that
\begin{equation}\label{bounds on sums}
    \sum_{n \in \II_l} \nu_n  \geq \frac {  2N+1  }2 \WW_0 (1-\eps).
%   \sum \nu_n  \geq \frac { (1- 2 \kappa)\,  | \II | }2 \WW_l (1-\eps)
% \quad \sum \si_n \leq {|\JJ|}\, \WW_1,
\end{equation}
\comm{****  since:

\ssk\nin (i)
 All horizontal arcs of $\AAA_l $ get pulled off at moments $n\in\II$;

\ssk\nin (ii)
The total weight of these arcs is at least $\WW_l (1-\eps)$;

\ssk\nin (iii)
Each of the arcs $(f^{  n(\alpha) -1  })^* \alpha$, where $\alpha\in \AAA_l^\ess$,
fits well  into $(1-2 \kappa)\,  |\II|$ domains
$(f^{m-1})^* \Om$ (Lemma \ref{fitting of ess  arcs});
% since the blocks $\BB_m$ with $|m|\leq (1-\kappa) N$  do not get lost
% in  the pullback process;

\ssk\nin (iv)
Each of these arcs has two ends: this yields  $1/2$;

\ssk\nin (v)
  At the moments $n\not\in \II$,  we do not gain any  horizontal weight.   

\bignote{More details !} 

*********} 
%
Let $K_m$, $m\in \II$,  be the family  of little Julia sets in $V$
(introduced in the beginning of the proof) 
and   let $\AAA_0$ be the set of horizontal arcs in $\BV$  that begin on
$K_0$ and fit into $ V$.  
 By Lemma  \ref{interesting moments},
 for any  horizontal arc $\beta \in \AAA_0$,  
the  pulled-off  moment $n(\beta) $ belongs to $\II$. 
Since the correspondence $\beta\mapsto (f^{n(\beta)-1} )^*\beta$
is at most  two-to-one, the total weight of the arcs  
$(f^{n(\beta)-1} )^*\beta$, $\beta\in \AAA_0$,  is at least $\WW_0/2$. 
Moreover,
% by Lemma \ref{fitting of arcs},  
each of the arcs $(f^{n(\beta)-1} )^*\beta$ fits into one of the $2N+1$ \note{explain}
domains $ (f^{m-1})^* \Om ' $, $m \in \II_0$,   \note{ $\Om'= (f^p)^* (\Om)$ } 
implying (\ref{bounds on sums}).

\bignote{In fact
  all the events in question happen within a bounded range of moments.
  Indeed, the loss of the short arc must happen at the Julia
  set $K_i$, $-\bb\leq i <0$  (before the short arc enters the
  translation region). At this moment one can form some arcs from $K_0$
  to $K_i$. These arcs must be lost at some $K_i$, $i\geq -2\bb$, etc.
  Hence the essential weight emanating from $K_0$ must be formed in bounded
  number of steps. }


Summing (\ref{loss at one step} ) up over $n\in \II$, taking into
account  (\ref{penalty}) and  (\ref{bounds on sums}), 
we see that the total horizontal loss is at least
\begin{equation}\label{loss of ess weight}
   \frac {2N+1}  4 [  (1-\eps)\WW_0 -\WW_1] - O(N^2 ).
%  \frac {|\II|} 4 [ (1- 2 \kappa) (1-\eps)\WW_l -\WW_1] - O(N^2 ).
% = \frac  {|\JJ|} 4  (\WW_0 -\WW_1 - O(N))
\end{equation}
So, if $\WW_0$ is definitely bigger than $\WW_1$, we are done. 

\comm{*******
\msk
Assume this is not the case.
% (Then (\ref{loss of ess weight}) shows only  that the gain of horizontal is
% bounded by $ N \eps'  \WW_\loc$.)  
Assume also that  $\WW_0$ is definitely bigger than
$\WW_1$ (and hence is definitely bigger than $\WW_l$  as well):
$$
     \WW_0- \WW_l \geq 2\de \WW_\loc,
$$
with an absolute $\de>0$. %  (to be selected below). 

By Lemma \ref{localization of local weight}, \note{!!} 
the total weight of  $\AAA_0$
is at least $(1-\eps)\WW_0$, where $\eps = O(1/\sqrt{\kappa N})$, so
it is definitely bigger than $\WW_l$.  Let $\tl \AAA_0$ be the family of those arcs
$\beta\in \AAA_0$ that are not  $f^l$-pullbacks of arcs $\alpha\in \AAA_l $.  
\bignote{It can be old arcs with a bigger weight -- correct counting (?) }
The
total weight of these arcs is at least 
$$
  (1-\eps)\WW_0 - \WW_l\geq \de \WW_\loc
$$ 
(if $\eps$ is small enough).
Each arc $\beta$ gets pulled off under $(f^n)^*$ at some moment $n(\beta)<p$.
In fact, $n(\beta)\leq p - \kappa \qq$ (if $\beta$ is essential),
for otherwise $(f^{n(\beta)-1} )^*\beta$
would be a short arc in the translation region (by Corollary \ref{loc weights comparison}),
which cannot be pulled off.
********}


 
\msk  
On the other hand, if $\WW_0$ is approximately equal to $\WW_1$, 
we replace estimate (\ref{bound on nu}) with 
$$
      (\WW_0-\WW_1) +  \nu_n/ 2 = \eps' \WW_\loc +\nu_n/2,
$$
which results in the  horizontal loss $\nu_n/2- \eps' \WW_\loc$.
Summing these up, taking (\ref{bounds on sums}) into account,
we produce, once again,  a total horizontal loss of order $N\WW_\loc$. 

% The rest of the argument is the same as in the bounded case for big
% enough $N$, 
% provided  $\kappa=\kappa(d)$ and $\eps=\eps(N)$  are sufficiently
% small while $\WW_\loc=\WW_\loc(N) $ is
%sufficiently big  (all independent of $p>>N$). 

                                                                                     